% Adaptive upper-tail (CVaR) advantages: findings so far
% Standalone article; the \section bodies can be pasted into report.tex frames as-is.
\documentclass[11pt]{article}
\usepackage[margin=1in]{geometry}
\usepackage{amsmath}
\usepackage{amssymb}
\usepackage{booktabs}
\usepackage{hyperref}

\title{Adaptive upper-tail (CVaR) advantages: findings so far}
\author{}
\date{4 September 2026}

\begin{document}
\maketitle

\noindent\emph{Based on the Erd\H{o}s runs under \texttt{runs/} (Qwen3-8B, 512 rollouts per step) and the analysis in \texttt{plot\_reward\_pdf.py} / \texttt{advantage.py}.}

\section{Setup}

The trainer's default objective is the entropic one from the paper. Per group of $K$ rollouts with rewards $R_1,\dots,R_K$, the advantages are
\[
  A_i = w_\beta(R_i) - 1, \qquad
  w_\beta(R_i) = \frac{e^{\beta R_i}}{\mathbb{E}\!\left[e^{\beta R}\right]},
\]
and $\beta$ is chosen per group by bisection so that the tilted distribution $q_\beta(i)\propto e^{\beta R_i}$ satisfies $\mathrm{KL}(q_\beta \,\|\, \mathrm{Uniform}) = \log 2$. As $\beta \to \infty$ this concentrates on the single best rollout.

\texttt{train\_multy\_CVaR.py} adds two alternatives selected by \texttt{--advantage-mode}. Plain GRPO uses $A_i = (R_i - \bar R)/(\sigma_R + \delta)$. The \texttt{cvar} mode blends GRPO with a sparse upper-tail term,
\[
  A_i = (1-\lambda)\,\frac{R_i - \bar R}{\sigma_R + \delta} + \lambda\, A_i^{\mathrm{Upper}}, \qquad
  A_i^{\mathrm{Upper}} =
  \begin{cases}
    \dfrac{R_i - q_{1-\alpha}}{\sigma_R + \delta} & R_i > q_{1-\alpha},\\[6pt]
    0 & \text{otherwise},
  \end{cases}
\]
with $q_{1-\alpha} = \mathrm{Quantile}(\mathbf R, 1-\alpha)$. The open question is how to set $\alpha$ per step instead of fixing it.



\section{Observation 1: the reward distribution collapses to a few discrete values}

Valid rewards go from 134 distinct values at step 0 to 37 by step 20, six by step 26 (98\,\% on one value) and a single value by step 33 in \texttt{erdos\_Qwen3-8B\_0827-1532}: the model replays a handful of constructions. With 30--75\,\% of every batch invalid at reward 0, each step is a spike at zero plus a tie cluster near the maximum, and $R = 1/c_5$ pushes all late-run progress into the fifth decimal of the reward.

\section{Observation 2: EVT threshold selection is unstable on this data}

The GPD-fit-plus-KS threshold rule (Clauset--Shalizi--Newman) is verified on synthetic data but erratic on real steps: $k^\ast$ jumps between 19, 96, 150, 25 and 134 on consecutive steps, most KS distances fail the goodness-of-fit test, and fitted shapes reach $-14$ with scales of $10^{-8}$. A continuous tail model cannot fit a staircase of ties on a $10^{-6}$ scale; fully tied steps have no excesses at all and fall back to ``tail $=$ ties at the maximum''.

\section{Observation 3: the Gini-modulated \texorpdfstring{$\alpha$}{alpha} is stable but uninformative}

Over all rewards, Gini is dominated by the zeros and equals roughly $1 - (\text{valid fraction})$, so $\alpha$ tracks validity, not tail shape. Over positive rewards Gini is $0.00$--$0.07$ every step, so $\alpha$ sits at $\alpha_{\mathrm{base}}$ and the rule is a fixed quantile.

\section{Observation 4: thresholds and ties}

Because thresholds are observed values and values are heavily tied, ``strictly above the threshold'' and ``top $k$'' disagree, sometimes reporting 0 rollouts above a threshold that 60 rollouts sit on. Any CVaR rule here must state that boundary ties belong to the tail; otherwise the tail is empty once the run converges.

\section{Implication}

The estimators are sound; the space they run in is not. The viable directions are, in order: analyse the log gap to target $\log(c_5 - c_5^{\mathrm{target}})$ instead of the reward, treat rewards as discrete levels and take the top $m$ levels as the tail, or gate the GPD fit on the KS test and smooth $\alpha_t$ across steps. Fully tied steps should be reported as having no tail rather than assigned an $\alpha$.

\end{document}
